\section{Análisis}

\subsection{Entregables del reto (Sección IV, documento de evaluación)}
\begin{enumerate}
  \item Gramática correcta y no ambigua para kernels Triton, con explicación de decisiones.
  \item Gestión de tablas de símbolos durante el análisis sintáctico.
  \item Implementación del parser con yacc.
  \item Mensajes de error sintácticos.
  \item Ejemplos de salida del parser.
\end{enumerate}

\subsection{Requerimientos funcionales}
\begin{itemize}
  \item Reconocer programas \texttt{.tri}: imports, \texttt{@triton.jit}, \texttt{def}, cuerpos indentados.
  \item Sentencias: asignación (\texttt{ident\_stmt}, anotada \texttt{id : tipo =}), \texttt{if/elif/else}, \texttt{while}, \texttt{for/in}, \texttt{return}, \texttt{pass}.
  \item Expresiones: literales, operadores (500), operador \texttt{is}, expresión condicional (\texttt{expr if expr else expr}), llamadas \texttt{tl.*}, subíndices \texttt{[:, None]}, argumentos con nombre (\texttt{mask=}, \texttt{dtype=}).
  \item Validación de línea: rechazo de múltiples sentencias en la misma línea (lexer + \texttt{validate\_stmt\_line\_end}).
  \item Indentación: \texttt{DEDENT} automático cuando una línea comienza en columna~0 tras un bloque indentado.
  \item Rechazar sintaxis inválida con \texttt{SYNTAX\_ERROR} en español.
  \item Actualizar symtab: imports, funciones, parámetros, variables.
\end{itemize}

\subsection{CFG (resumen)}
La gramática formal completa (\textbf{36 no terminales}, \textbf{105 alternativas} numeradas P-001\ldots P-105, inventario en Tabla~\ref{tab:inventario-nt} del Apéndice~\ref{ap:cfg-completa}) se documenta en el Listado~\ref{lst:cfg-completa} y el bloque yacc Listado~\ref{lst:parser-y-gramatica}. El siguiente extracto resume las construcciones principales para lectura rápida en el análisis.

\begin{lstlisting}[caption={CFG resumida del lenguaje .tri (Triton kernel subset)}]
program       -> opt_newlines top_block opt_newlines
top_block     -> top_block top_line | epsilon
top_line      -> stmt | stmt NEWLINE

stmt          -> import_stmt | def_stmt | assign_stmt | ident_stmt | if_stmt
              | while_stmt | for_stmt | return_stmt | pass_stmt | expr_stmt

ident_stmt    -> TOKEN_IDENTIFIER ident_op
ident_op      -> TOKEN_ASSIGN expr | ident_suffix

block_stmt    -> stmt TOKEN_NEWLINE | stmt   /* validate_stmt_line_end */

import_stmt   -> KW_IMPORT import_name opt_as
import_name   -> ID | import_name '.' ID

opt_triton_decorator -> epsilon | '@' ID '.' ID opt_newlines

def_stmt      -> opt_triton_decorator KW_DEF ID '(' param_list_opt ')' ':' suite

suite         -> simple_stmt
              | opt_newlines INDENT block_body DEDENT

expr          -> literal | ID ident_suffix
              | expr OP expr | expr KW_IS expr
              | expr KW_IF expr KW_ELSE expr
              | expr '[' subscript_list ']'
              | '(' arg_list_opt ')' | call_expr | unary_expr

subscript     -> expr | NONE | ':' opt_expr | opt_expr ':' opt_expr
              | opt_expr ':' opt_expr ':' opt_expr

call_expr     -> ID '(' arg_list_opt ')'
              | ID member_chain '(' arg_list_opt ')'

arg           -> expr | ID '=' expr
\end{lstlisting}

\subsection{Pasos para obtener la gramática}
\begin{enumerate}
  \item Inventario de tokens del reporte léxico (Cuadro~1).
  \item Identificación de estructuras Triton en 70 kernels de referencia.
  \item Borrador de producciones para bloques indentados (\texttt{INDENT}/\texttt{DEDENT}).
  \item Incorporación de expresiones, llamadas y subíndices tensoriales.
  \item Resolución de ambigüedades con precedencias yacc y pruebas iterativas.
\end{enumerate}

\subsection{Simplificaciones justificadas}
\begin{table}[H]
\centering
\caption{Simplificaciones de la CFG.}
\begin{tabular}{@{}p{3.2cm}p{8.5cm}@{}}
\toprule
Tipo & Decisión \\
\midrule
Unitarias & Jerarquía de expresiones colapsada en \texttt{expr} con \texttt{\%left}. \\
$\varepsilon$ & Solo en listas opcionales, \texttt{else} y líneas vacías iniciales. \\
Símbolos inútiles & Sin \texttt{class}, \texttt{try}, \texttt{switch} (no usados en kernels). \\
Indentación & \texttt{WS\_COUNT} $\rightarrow$ pila $\rightarrow$ \texttt{INDENT}/\texttt{DEDENT}. \\
Comentarios & Patrón \texttt{\#[\^{})\textbackslash n]*} para no consumir \texttt{)} en la misma línea. \\
\bottomrule
\end{tabular}
\end{table}

\subsection{Modificaciones al analizador léxico}
\begin{itemize}
  \item Integración con \texttt{y.tab.h}; \texttt{yylex()} retorna códigos de token.
  \item Tokens sintéticos 910/911 (no en Cuadro~1 original; documentados en este reporte).
  \item Sin nuevos IDs en el Cuadro~1 oficial; trazabilidad numérica conservada.
  \item \textbf{Detección \texttt{BAD\_ID}} (heredada del analizador léxico del equipo en \url{https://github.com/G4ndhiV/Analizador-lexico}): identificadores que contienen caracteres prohibidos (\texttt{\$}, acento grave, barra invertida, \texttt{?}, \texttt{!}) se reconocen con una regla Flex dedicada \emph{antes} que \texttt{ID}, emitiendo \texttt{LEXICAL\_ERROR} (999) sin fragmentar el lexema.
  \item \textbf{Múltiples sentencias en la misma línea}: el wrapper \texttt{yylex()} detecta patrones inválidos (\texttt{pass pass}, \texttt{return 5 6}, \texttt{x = 5 y = 6}, etc.) y marca \texttt{syntax\_error\_seen}.
  \item \textbf{DEDENT en columna~0}: si un identificador inicia una línea sin espacios previos y la pila de indentación está activa, se emite \texttt{TOKEN\_DEDENT} antes del token (cierre de bloque Python).
\end{itemize}

\subsection{Patrones léxicos: \texttt{BAD\_ID}}
\begin{table}[H]
\centering
\caption{Caracteres invalidos en identificadores y patron Flex.}
\label{tab:bad-id}
\begin{tabular}{@{}llp{7.2cm}@{}}
\toprule
\textbf{Definición} & \textbf{Regex} & \textbf{Comportamiento} \\
\midrule
\texttt{BAD\_ID\_CHAR} & \texttt{[\$`{\textbackslash}?!]} & Conjunto de símbolos no permitidos en nombres. \\
\texttt{BAD\_ID} & ver \texttt{lexer/lexico\_equipo\_11\_scanner.l} & Secuencias con al menos un carácter inválido; prioridad sobre \texttt{ID}. \\
\texttt{ID} & \texttt{[a-zA-Z\_][a-zA-Z0-9\_]*} & Identificadores válidos (tabla 100). \\
\bottomrule
\end{tabular}
\end{table}
Mensaje emitido: \texttt{LEXICAL\_ERROR: ... (identificador con caracteres invalidos)}. El driver marca \texttt{lexical\_error\_seen} y termina con código de salida~2, sin reportar \texttt{PARSE\_OK}.

\subsection{Trazabilidad token léxico $\rightarrow$ yacc}
\begin{table}[H]
\centering
\caption{Correspondencia Token ID (reporte léxico) y yacc.}
\begin{tabular}{@{}cll@{}}
\toprule
ID & Nombre & yacc \\
\midrule
100 & IDENTIFIER & \texttt{TOKEN\_IDENTIFIER} \\
101 & KEYWORD & \texttt{KW\_def}, \texttt{KW\_if}, \ldots \\
200--400 & Literales & \texttt{TOKEN\_INTEGER}, etc. \\
500 & OPERATOR & \texttt{TOKEN\_OPERATOR} \\
600 & DELIMITER & \texttt{TOKEN\_DELIMITER} o literales \texttt{( ) [ ] : ,} \\
900 & WS\_COUNT & convertido a INDENT/DEDENT \\
950 & NEWLINE & \texttt{TOKEN\_NEWLINE} \\
910/911 & (sintéticos) & \texttt{TOKEN\_INDENT}, \texttt{TOKEN\_DEDENT} \\
\bottomrule
\end{tabular}
\end{table}
